% ============================================================================
%  Guion de defensa — TP Final SdS (72.25) — ITBA — 2026 Q1 — Grupo G01S2
%  Sincronizado con SdS_TPFinal_2026Q1G01S2_Presentacion.pdf (18 diapositivas).
%  Compilar: pdflatex guion-defensa_v1.tex
% ============================================================================
\documentclass[11pt,a4paper]{article}

\usepackage[utf8]{inputenc}
\usepackage[T1]{fontenc}
\usepackage{lmodern}
\usepackage{microtype}
\usepackage[margin=2.2cm]{geometry}
\usepackage{xcolor}
\usepackage{booktabs}
\usepackage{amssymb}
\usepackage{enumitem}
\usepackage[hidelinks]{hyperref}
\usepackage{parskip}

% Colores por orador (para ubicar rápido la parte de cada uno al leer online)
\definecolor{cnico}{RGB}{20,90,170}
\definecolor{cian}{RGB}{175,45,45}
\definecolor{ckeo}{RGB}{30,130,70}

% Encabezado de cada diapositiva: nro, titulo, orador (color), tiempo
\newcommand{\diapo}[5]{%
  \medskip\par\noindent
  \textcolor{#4}{\textbf{$\blacktriangleright$\ Diapo #1 --- #2}}\quad
  \textcolor{#4}{\textbf{[#3]}}\hfill{\itshape\small(#5)}\par\nobreak\smallskip
}
\newcommand{\relevo}[1]{\hfill\textbf{\itshape #1}}
\newcommand{\videocue}[1]{\textcolor{cian!85!black}{\textbf{[$\blacktriangleright$ REPRODUCIR: #1]}}\ }

\setlength{\parindent}{0pt}

\begin{document}

\begin{center}
  {\LARGE\textbf{Guion de defensa --- TP Final SdS}}\\[4pt]
  {\large Diagrama fundamental de vehículos dirigidos por vibración (VDV)}\\[6pt]
  72.25 Simulación de Sistemas --- ITBA --- 2026 Q1 --- \textbf{Grupo G01S2}\\[4pt]
  \textcolor{cnico}{\textbf{Nico}} \; · \; \textcolor{cian}{\textbf{Ian}} \; · \; \textcolor{ckeo}{\textbf{Keo}}
  \hspace{1.5em}|\hspace{1.5em}\textbf{Objetivo: $\sim$12 min}\\[2pt]
  {\small Sincronizado con las 18 diapositivas de \texttt{Presentacion.pdf}. Online: cada uno habla su bloque de corrido.}
\end{center}

\vspace{2pt}
\hrule
\vspace{6pt}

\section*{Reparto (tres bloques de corrido, solo dos relevos)}

\begin{center}
\begin{tabular}{llll}
\toprule
\textbf{Bloque} & \textbf{Diapos} & \textbf{Presenta} & \textbf{Tiempo} \\
\midrule
Contexto y modelo & 1--6 & \textcolor{cnico}{\textbf{Nico}} & $\sim$4:00 \\
Motor, corridas y dinámica ($N$ fijo) & 7--12 & \textcolor{cian}{\textbf{Ian}} & $\sim$4:15 \\
Resultados del orden, validación y cierre & 13--18 & \textcolor{ckeo}{\textbf{Keo}} & $\sim$4:05 \\
\bottomrule
\end{tabular}
\end{center}

\textbf{Videos:} las diapos \textbf{9} y \textbf{12} tienen animaciones (los fotogramas son clicables y abren
el \texttt{.mp4}). Si el visor bloquea el click, tengan los \texttt{.mp4} de \texttt{data\_anim/} abiertos en un
reproductor y los pasan al compartir pantalla. Dejen correr el video mientras leen el texto de esa diapo.

\textbf{Tips:} ritmo tranquilo ($\sim$150 palabras/min). Tomen el tiempo al ensayar y recorten donde se estire.

\vspace{4pt}
\hrule
\vspace{4pt}

\section*{Guion}

% ============================ BLOQUE 1 --- NICO (1-6) ========================
\diapo{1}{Portada}{NICO}{cnico}{20 s}
Hola, buenas. Somos el grupo G01S2: Nico Arias, Keo Dubovitsky e Ian Tognetti. Lo que estudiamos es el
\textbf{diagrama fundamental de vehículos dirigidos por vibración}: agarramos un autómata celular de
Nagel--Schreckenberg, le metimos una interacción por contacto, y con eso reproducimos las \textbf{tendencias}
que Patterson y Parisi midieron en el experimento. Arranco con el sistema real.

\diapo{2}{El sistema real: VDV}{NICO}{cnico}{45 s}
El sistema son unos robots comerciales, los Hexbug Nano. Se mueven porque la vibración se rectifica sobre unas
cerditas asimétricas. Los meten en un canal \textbf{circular} de unos 1313 milímetros, así que es un sistema de
\textbf{una dimensión}, y entran hasta 30. Lo clave ---y por eso se usa este sistema--- es que
\textbf{interactúan solo por contacto}: no tienen forma de ``ver'' al de adelante para frenar, como sí hace un
auto. Cada uno tiene su velocidad libre, repartidas más o menos parejo entre 90 y 120 milímetros por segundo.
Con esto se \textbf{aísla} la parte del diagrama fundamental que sale solo de los choques físicos.

\diapo{3}{Lo que mide el experimento}{NICO}{cnico}{45 s}
¿Qué es lo que mide el experimento? Cuatro cosas. Una: a densidad baja y media la velocidad se queda
\textbf{casi constante}. Dos: \textbf{cae entre un 25 y un 40 por ciento} cuando te acercás a la densidad de
contacto. Tres: importa el \textbf{orden} en que vas metiendo los vehículos ---el aleatorio queda \textbf{acotado}
entre el creciente y el decreciente---. Y cuatro: cuando llegás a los 30, las tres configuraciones
\textbf{convergen}, a una velocidad que encima queda por debajo de la del más lento. Nuestra idea es reproducir
estas \textbf{tendencias} con un modelo barato.

\diapo{4}{Modelo: NaSch con contacto}{NICO}{cnico}{50 s}
El modelo es un autómata celular de Nagel--Schreckenberg en una \textbf{ruta periódica de una dimensión}. Cada
vehículo ocupa $\ell$ celdas y su velocidad es entera. La actualización es \textbf{sincrónica}, en el orden R1,
R3, R2, R4: acelerar, frenar al azar, resolver los contactos, y mover. Lo importante ---y lo hicimos a
propósito--- es que \textbf{frenamos antes de resolver los contactos}: como frenar solo achica el avance,
resolver después nos garantiza que nunca se solapen. Lo que le cambiamos es la Regla 2, la interacción: la
hicimos de \textbf{contacto puro}. El vehículo avanza libre hasta que alcanza al de adelante; si lo alcanzaría,
queda pegadito a contacto y le \textbf{hereda la velocidad}. Esa fórmula del desplazamiento asegura que no se
encimen.

\diapo{5}{Dos variantes de la Regla 2}{NICO}{cnico}{35 s}
Tenemos dos variantes de esa Regla 2. La \textbf{oficial, la A}, es contacto puro: sin anticipar nada, flujo
libre hasta el contacto, y al alcanzar al de adelante le hereda la velocidad. Esa es la física del experimento.
La otra, la \textbf{B}, es la clásica de manual, que usamos \textbf{solo para validar} el motor contra la
solución analítica. Y ojo: la de contacto puro \textbf{no} se valida contra la curva clásica, son modelos
distintos.

\diapo{6}{Calibración}{NICO}{cnico}{30 s}
Y calibramos todo a la geometría del experimento: el paso espacial es un cuarto de milímetro; el temporal, uno
sobre veinticuatro de segundo ---24 cuadros por segundo, como la cámara---; el vehículo mide 44 milímetros; y la
pista la elegimos de 1320, que es justo 30 veces el largo del vehículo. Con eso, los 30 vehículos llenan la ruta
exactamente a contacto. \relevo{Ahí te paso, Ian, con la implementación.}

% ============================ BLOQUE 2 --- IAN (7-12) =======================
\diapo{7}{Implementación del motor}{IAN}{cian}{35 s}
Gracias. El motor está hecho en Java. Cada paso se calcula sobre una \textbf{instantánea inmutable} de los
huecos y las velocidades, así que el resultado \textbf{no depende del orden} en que recorrés los vehículos. El
no-solapamiento está garantizado por construcción, y la inserción incremental la hacemos \textbf{por empuje}, que
nos deja llegar a los 30 exactos. Y algo que la cátedra nos remarca: el motor es \textbf{determinista} y saca
\textbf{solo variables físicas}, posición y velocidad; el color lo calculamos después.

\diapo{8}{Simulaciones y observables}{IAN}{cian}{40 s}
Del lado de las corridas: barremos el número de vehículos de 5 a 30, y la probabilidad de frenado de 0 a 0,4.
Usamos dos protocolos, uno de $N$ fijo y el incremental que va de 5 a 30, con \textbf{30 realizaciones} por punto.
Y algo clave para la cátedra: todos los observables se calculan \textbf{después} de la simulación, nunca adentro
del motor. Sacamos la velocidad media ---con el error como el \textbf{desvío entre realizaciones}---, la densidad
de cada uno, las distribuciones y el diagrama fundamental. El estacionario lo elegimos \textbf{a ojo, por
inspección}.

\diapo{9}{Dinámica $N$ fijo --- ANIMACIONES}{IAN}{cian}{55 s}
\videocue{N=5 y N=30 (izq.\ y der.)}Acá van las primeras animaciones, con $N$ fijo. Le doy play a la de la
izquierda: \textbf{$N=5$}, densidad baja. Miren que los vehículos van casi sin tocarse, cada uno a lo suyo ---el
\textbf{color es la velocidad}, verde es rápido---. Ahora la de la derecha: \textbf{$N=30$}, la ruta llena. Se ve
el anillo rígido, todos pegados a contacto, moviéndose juntos. Lo que \textbf{cambia} entre las dos es la
densidad, o sea cuántos vehículos hay; y esto lo vamos a \textbf{resumir con la velocidad media}, que sale en las
próximas dos diapos.

\diapo{10}{Estacionario por inspección}{IAN}{cian}{30 s}
Antes de medir, elegimos el corte del \textbf{estacionario}, y lo hacemos \textbf{a ojo}, mirando la serie
temporal de la velocidad media ---no descartando un porcentaje fijo, que es lo que pide la cátedra---. Se ve que,
pasado el transitorio, la serie queda fluctuando alrededor de un valor estable, y de ahí para adelante
promediamos.

\diapo{11}{Velocidad media vs $N$ ($N$ fijo)}{IAN}{cian}{40 s}
Y esta es la curva respuesta--estímulo con $N$ fijo: velocidad media contra $N$, una curva por cada $p$. A
densidad baja se queda casi constante, pero ojo con el detalle: esa meseta \textbf{no} está en la velocidad libre
media, que es 105, sino en la del \textbf{más lento}, ahí en 90. Es que en un solo carril, sin poder pasar, todo
se amontona atrás del más lento. Y en $N=30$, la ruta justo llena, cuánto baja depende de $p$.

\diapo{12}{Efecto del orden --- ANIMACIONES}{IAN}{cian}{55 s}
\videocue{los 3 órdenes (creciente / decreciente / aleatorio)}Ahora las animaciones del protocolo incremental, que
arranca en 5 y va sumando de a 5 hasta 30. Las tres son al mismo $p$ y muestran los tres órdenes: a la izquierda
\textbf{creciente} ---primero los lentos---, en el medio \textbf{decreciente} ---primero los rápidos---, y a la
derecha \textbf{aleatorio}. Le doy play... Fíjense que a \textbf{igual $N$} el color, o sea la velocidad, cambia
según el orden: el decreciente arranca más verde, más rápido, y el creciente más lento. Lo que \textbf{cambia}
acá es el \textbf{orden de inserción}, y de nuevo lo vamos a resumir con la velocidad media. \relevo{Ahí te paso,
Keo, con la curva.}

% ============================ BLOQUE 3 --- KEO (13-18) ======================
\diapo{13}{Velocidad media vs $N$ por orden}{KEO}{ckeo}{45 s}
Gracias, Ian. Efectivamente, esta es la curva análoga a la \textbf{Figura 2} del paper: la velocidad media según
el orden. El \textbf{decreciente} da la más alta, el \textbf{creciente} la más baja, y el \textbf{aleatorio}
queda acotado en el medio ---justo lo que pasa en el experimento---. Y las tres terminan \textbf{convergiendo a
saturación}, ahí en 81 milímetros por segundo con $p=0{,}1$, que ya está por debajo de los 90 del más lento.

\diapo{14}{Distribuciones por orden}{KEO}{ckeo}{40 s}
Acá están las distribuciones por $N$ activo, análogas a las \textbf{Figuras 3 y 4}. A la izquierda, la
\textbf{densidad}: el pico cae en la densidad de contacto, 1 sobre 44, y se hace más marcado a medida que sube
$N$. A la derecha, la \textbf{velocidad}, para el orden creciente: es casi independiente de $N$, porque en el
creciente el más lento ya está desde el arranque. En el decreciente, en cambio, la distribución se corre para
abajo a medida que van entrando los más lentos.

\diapo{15}{Diagrama fundamental por orden}{KEO}{ckeo}{35 s}
El \textbf{diagrama fundamental} por orden confirma toda la estructura: el decreciente es la curva de arriba, el
creciente la de abajo, y el aleatorio queda acotado entre las dos. Y la velocidad se cae al acercarse a la
densidad de contacto, igual que en el experimento.

\diapo{16}{Validación del motor}{KEO}{ckeo}{40 s}
Antes de las limitaciones, la \textbf{validación}. Con la variante clásica, la B, en su régimen ---partículas
puntuales, homogénea, $p=0$--- hay solución analítica: el diagrama \textbf{triangular}. Y el motor lo reproduce a
\textbf{precisión de máquina}, con tests automáticos a tolerancia $10^{-9}$. De nuevo el punto: esto valida
\textbf{el motor}; el contacto puro no se valida contra esta curva.

\diapo{17}{Limitaciones}{KEO}{ckeo}{45 s}
Y lo más \textbf{honesto}, las limitaciones. Primera: como imponemos no-solapamiento duro, \textbf{no}
reproducimos la cola de densidad por arriba del contacto que sí ve el experimento. Segunda: el punto $N=30$ es
\textbf{singular} y depende de que fijemos $L$ en 1320; el experimento tiene 1313. Y tercera: la caída por debajo
del más lento necesita $p>0$, y el mecanismo es \textbf{distinto} al del experimento ---de hecho, con $p\ge0{,}3$
el anillo lleno se \textbf{congela}, que ya no es ningún régimen real---; por eso el punto que sí es comparable es
$p=0{,}1$.

\diapo{18}{Conclusiones}{KEO}{ckeo}{40 s}
Y para cerrar. El autómata con \textbf{contacto puro} reproduce las \textbf{tendencias} centrales del diagrama
fundamental: velocidad casi plana a baja densidad y caída cerca del contacto. El \textbf{orden de inserción}
modula el diagrama, con el aleatorio acotado entre el creciente y el decreciente. Reproducimos el pico de
densidad en el contacto y el angostamiento de la velocidad con $N$. Y la variante clásica \textbf{valida el
motor} contra la curva triangular. Bueno, muchas gracias, y quedamos para las preguntas que tengan.

\vspace{6pt}
\hrule
\vspace{4pt}

\section*{Machete de preguntas típicas del jurado}

\begin{itemize}[leftmargin=1.2em,itemsep=3pt]
\item \textbf{¿Por qué $L=1320$ y no 1313 del paper?} Para que $N=30$ cierre justo a contacto ($30\,\ell$);
  el 1313 lo discutimos como sensibilidad. Con $L=1313$, $N=30$ no entra (5252 $<$ 5280 celdas).
\item \textbf{¿Qué valida exactamente la variante B?} El diagrama fundamental triangular analítico
  $Q(\rho)=\min(\rho\,v_{\max},\,1-\rho)$, en régimen homogéneo $\ell=1$, $p=0$. El contacto puro NO se
  valida contra esa curva.
\item \textbf{¿Cómo eligieron el estacionario?} A ojo, por inspección de la evolución temporal (no un
  porcentaje fijo); el corte queda registrado en \texttt{figures/manifiesto.csv}.
\item \textbf{¿Por qué la meseta está en $\sim$90 y no en 105?} Pista de un solo carril, sin adelantamiento:
  todo se amontona atrás del más lento, así que la media queda anclada a su velocidad, no a la libre media.
\item \textbf{¿El error?} Desvío \emph{entre} realizaciones (ddof$=1$), no el de todos los cuadros ni SEM.
\item \textbf{¿Qué NO reproduce el modelo?} La cola de densidad arriba del contacto, la forma de la
  distribución de velocidad (pico en $\sim$90 contra $\sim$65 del experimento) y el mecanismo del colapso
  (nuestro colapso con $p$ alto es un congelamiento artificial del anillo lleno).
\end{itemize}

\end{document}
